\section{Explicit Constructions of Expanders}
\label{sec:constructions}

The existence of expander families follows easily from the probabilistic method, since a random \(d\)-regular graph is an expander with high probability for any fixed \(d\geq 3\). The deeper and more demanding question is whether expander families can be constructed explicitly, meaning whether there exists a polynomial-time algorithm that, given \(n\), outputs a \(d\)-regular expander on \(n\) vertices with \(\SpectralGap\geq c>0\) for absolute constants \(d\) and \(c\). Explicit constructions are essential for all algorithmic and cryptographic applications, since one must be able to evaluate the graph efficiently.

\subsection{Probabilistic Existence}
\label{subsec:probabilistic-existence}

Before discussing explicit constructions, it is useful to isolate the simple probabilistic reason that expanders should exist in abundance. The guiding intuition is that a random bounded-degree graph has no reason to respect any proposed cut. If a set \(S\) of size \(r\) had only \(\bigO(r)\) external neighbors, then all edges leaving \(S\) would have to land inside a prescribed set \(S'\) of size only slightly larger than \(S\), and the probability of this happening decays exponentially in the number of random edges exposed from \(S\).

The random-permutation model below is deliberately elementary. It avoids the dependencies present in the uniform simple regular model, while still producing a bounded-degree graph in which every left endpoint has a fixed number of random images, and this is enough to show that expansion is not a rare phenomenon.

\begin{theorem}[Random Permutation Multigraphs Expand]
	\label{thm:random-permutation-expanders}
	There are constants \(\ell_0\in\NN\) and \(\alpha>0\) such that the following holds for every fixed \(\ell\geq \ell_0\). Let \(\pi_1,\ldots,\pi_\ell\) be independent uniformly random permutations of \([n]\), and let \(\Graph_n\) be the undirected \(2\ell\)-regular multigraph obtained by adding, for every \(i\in[\ell]\) and every \(v\in[n]\), the edge \(\{v,\pi_i(v)\}\), where loops and repeated edges are counted with multiplicity in the degree. Then
	\begin{equation}
		\label{eq:random-permutation-expansion-probability}
		\Pr[\mathsf{Exp}_n(\alpha)]=1-o_{n\to\infty;\ell}(1),
	\end{equation}
	where \(\mathsf{Exp}_n(\alpha)\) denotes the event that \(\Graph_n\) is an \(\alpha\)-vertex expander.
\end{theorem}

\begin{proof}
	The proof follows the standard first-moment strategy. We fix a possible bad set \(S\), enlarge it to the small set \(S'\) that would contain all of its neighbors if expansion failed, and bound the probability that every random permutation maps \(S\) into \(S'\). The only point that needs care is that we must union-bound over all choices of \(S\) and \(S'\), and the degree parameter \(\ell\) is chosen large enough to make the exponential probability loss dominate this entropy.

	Fix a constant \(c\in(0,1/10)\), to be chosen small enough below, and put
	\begin{equation}
		\label{eq:random-permutation-r-prime}
		r'=\lfloor cr\rfloor+1
	\end{equation}
	for each \(1\leq r\leq n/2\). We shall prove that, with probability \(1-o_{n\to\infty;\ell}(1)\), there do not exist sets \(S\subseteq S'\subseteq[n]\) such that
	\begin{equation}
		\label{eq:random-permutation-bad-sizes}
		\abs{S}=r,\quad \abs{S'}=r+r',
	\end{equation}
	and
	\begin{equation}
		\label{eq:random-permutation-all-images-contained}
		\forall i\in[\ell],\quad \pi_i(S)\subseteq S'.
	\end{equation}
	Indeed, if some set \(S\) with \(1\leq\abs{S}\leq n/2\) had fewer than \(c\abs{S}\) external neighbors in \(\Graph_n\), then the set \(S'=S\cup\extNbr_{\Graph_n}(S)\), enlarged arbitrarily if necessary to have size \(\abs{S}+\lfloor c\abs{S}\rfloor+1\), would satisfy \cref{eq:random-permutation-bad-sizes,eq:random-permutation-all-images-contained}. Thus excluding such pairs \(S,S'\) proves \(\alpha\)-vertex expansion for any \(\alpha<c\), after decreasing \(\alpha\) to absorb the harmless rounding in \cref{eq:random-permutation-r-prime}.

	For fixed sets \(S\subseteq S'\) satisfying \cref{eq:random-permutation-bad-sizes}, the probability that one random permutation maps \(S\) into \(S'\) is
	\begin{equation}
		\label{eq:single-permutation-containment-probability}
		\Pr[\pi_i(S)\subseteq S']
		=
		\frac{\binom{r+r'}{r}}{\binom{n}{r}}
		\leq
		\left(\frac{r+r'}{n}\right)^r.
	\end{equation}
	The equality holds because the image \(\pi_i(S)\) is a uniformly random \(r\)-element subset of \([n]\), and the inequality follows by writing both binomial coefficients as products and comparing term by term. Since the permutations are independent, \cref{eq:single-permutation-containment-probability} gives
	\begin{equation}
		\label{eq:all-permutations-containment-probability}
		\Pr\left[\bigcap_{i\in[\ell]}\{\pi_i(S)\subseteq S'\}\right]
		\leq
		\left(\frac{r+r'}{n}\right)^{\ell r}.
	\end{equation}

	We now count the possible pairs \(S\subseteq S'\). There are
	\begin{equation}
		\label{eq:bad-pair-count-exact}
		\binom{n}{r}\binom{n-r}{r'}
		=
		\frac{n!}{r!\,r'!\,(n-r-r')!}
	\end{equation}
	such pairs. A crude Stirling bound gives an absolute constant \(C>1\) such that, for all \(r\leq n/2\),
	\begin{equation}
		\label{eq:bad-pair-count-stirling}
		\frac{n!}{r!\,r'!\,(n-r-r')!}
		\leq
		C^r\left(\frac{n}{r+r'}\right)^{r+r'}.
	\end{equation}
	For completeness, one may see \cref{eq:bad-pair-count-stirling} directly by writing \(a=r+r'\), observing that \(a\leq 2r+1\), and using
	\begin{equation}
		\label{eq:bad-pair-count-stirling-derivation}
		\binom{n}{r}\binom{n-r}{r'}
		=\binom{n}{a}\binom{a}{r}
		\leq \left(\frac{en}{a}\right)^a2^a
		\leq C^r\left(\frac{n}{a}\right)^a,
	\end{equation}
	after increasing \(C\) to handle the finitely many small values of \(r\).
	Combining \cref{eq:all-permutations-containment-probability,eq:bad-pair-count-stirling}, the probability that a bad pair exists for this value of \(r\) is at most
	\begin{equation}
		\label{eq:bad-pair-probability-r}
		\Pr[\mathsf{Bad}_r]\leq
		C^r
		\left(\frac{r+r'}{n}\right)^{\ell r-r-r'}.
	\end{equation}
	Here \(\mathsf{Bad}_r\) denotes the event that there exist sets \(S\subseteq S'\subseteq[n]\) satisfying \cref{eq:random-permutation-bad-sizes,eq:random-permutation-all-images-contained}.
	Because \(c<1/10\), \cref{eq:random-permutation-r-prime} implies \(r+r'\leq 2r\) for every \(r\geq 1\) and \(r+r'\leq 0.6n\) for every \(r\leq n/2\), once \(n\) is large. We split the summation over \(r\). If \(1\leq r\leq n^{1/2}\), then \cref{eq:bad-pair-probability-r} gives
	\begin{equation}
		\label{eq:random-permutation-small-r-bound}
		C^r
		\left(\frac{r+r'}{n}\right)^{\ell r-r-r'}
		\leq
		\left(C\left(\frac{2r}{n}\right)^{\ell-2}\right)^r
		\leq
		\left(C2^{\ell-2}n^{-(\ell-2)/2}\right)^r,
	\end{equation}
	and the sum of \cref{eq:random-permutation-small-r-bound} over this range is \(o_{n\to\infty;\ell}(1)\) for every fixed \(\ell\geq 3\). If \(n^{1/2}<r\leq n/2\), then choosing \(\ell_0\) so large that
	\begin{equation}
		\label{eq:random-permutation-degree-choice}
		C\cdot 0.6^{\ell_0-2}<\frac{1}{2}
	\end{equation}
	makes the contribution of each such \(r\) at most \(2^{-r}\). Therefore
	\begin{equation}
		\label{eq:random-permutation-union-bound}
		\Pr\left[\bigcup_{1\leq r\leq \lfloor n/2\rfloor}\mathsf{Bad}_r\right]
		\leq
		o_{n\to\infty;\ell}(1)+\sum_{r>n^{1/2}}2^{-r}
		=o_{n\to\infty;\ell}(1),
	\end{equation}
	as required. As argued above, the absence of bad pairs implies vertex expansion with a constant \(\alpha>0\) depending only on \(c\), and therefore only on the fixed choice of \(\ell_0\).
\end{proof}

\cref{thm:random-permutation-expanders} is intentionally stated for multigraphs, because this is the cleanest model for the union-bound argument. One may condition on the absence of loops and repeated edges, or work directly in the uniform simple \(d\)-regular model, to obtain the standard statement that random regular graphs are expanders with high probability; see Pinsker~\cite{Pin1973}, Bollobas~\cite{Bol1988}, and the survey of Hoory, Linial, and Wigderson~\cite{HLW2006}. The lesson for this note is conceptual rather than model-theoretic: randomness produces expansion easily, while explicitness requires additional algebraic or combinatorial structure.

\subsection{Algebraic Constructions of Ramanujan Graphs}
\label{subsec:ramanujan}

The probabilistic construction above explains why expanders exist, but it does not explain how to name one, evaluate its neighbor function, or place it inside an algorithm. The first successful explicit constructions therefore had to replace randomness by structure. Margulis~\cite{Mar1973} did this by using group actions, and the later construction of Lubotzky, Phillips, and Sarnak~\cite{LPS1988} showed that the same algebraic viewpoint can even meet the spectral limit imposed by the Alon--Boppana theorem in \cref{thm:alon-boppana}.

We say that a \(d\)-regular graph is Ramanujan if every nontrivial adjacency eigenvalue \(\theta\), excluding the eigenvalues forced by connected components and bipartiteness, satisfies
\begin{equation}
	\label{eq:ramanujan-bound}
	\abs{\theta}\leq 2\sqrt{d-1}.
\end{equation}
In normalized random-walk form, \cref{eq:ramanujan-bound} says that every nontrivial eigenvalue of \(\Markov{\Graph}\) has absolute value at most \(2\sqrt{d-1}/d\). Thus Ramanujan graphs are not merely expanders; they are spectrally optimal bounded-degree expanders, up to the unavoidable obstruction captured by \cref{thm:alon-boppana}.

The conceptual reason that algebra enters the construction is that Cayley graphs turn spectral analysis into representation theory. A Cayley graph has vertices given by a finite group and edges given by multiplication by a fixed generating set. Once the generating set is chosen symmetrically, the graph is regular and undirected, and the adjacency operator becomes an averaging operator over group elements. The problem of bounding eigenvalues is then converted into the problem of understanding how this averaging operator acts in every nontrivial representation of the group.

\begin{construction}[LPS Ramanujan Graphs~\cite{LPS1988}]
	\label{con:lps}
	Let \(p\) and \(q\) be distinct odd primes satisfying the congruence conditions required in the LPS construction. The graph \(X^{p,q}\) is a Cayley graph of either \(\mathrm{PGL}(2,\FF_q)\) or \(\mathrm{PSL}(2,\FF_q)\). Its \(p+1\) generators are obtained from integer solutions to
	\begin{equation}
		\label{eq:lps-four-square-equation}
		\alpha_0^2+\alpha_1^2+\alpha_2^2+\alpha_3^2=p
	\end{equation}
	with prescribed parity conditions. Multiplication by these generators defines a \((p+1)\)-regular graph. The resulting graph is Ramanujan, and in the \(\mathrm{PSL}(2,\FF_q)\) case the number of vertices is \(q(q^2-1)/2\). The proof of the Ramanujan bound uses deep number theory, ultimately relying on Deligne's proof of the relevant Ramanujan conjecture for modular forms.
\end{construction}

For the purposes of this note, the details hidden in \cref{con:lps} are less important than the structural lesson. Algebra can produce graphs whose local rule is explicit and whose global mixing is essentially as good as possible, but the proof of this fact may use machinery far beyond elementary graph theory. This is one reason that the combinatorial constructions below are so important, since they explain how expansion can be built and maintained by operations on graphs, rather than by appealing to deep eigenvalue estimates from number theory.

\subsection{Combinatorial Constructions via the Zig-Zag Product}
\label{subsec:zig-zag}

The zig-zag product of Reingold, Vadhan, and Wigderson~\cite{RVW2002} gives a different answer to the explicitness problem. Rather than constructing one highly structured graph from a group, it gives an operation that combines two expanders into a new expander while controlling the degree. The operation is best understood as solving a tension that appears whenever one tries to square a graph to improve expansion: graph squaring improves the spectral gap, but it also squares the degree. The zig-zag product restores the degree by replacing each high-degree vertex by a small expander, called a cloud, and using the small graph to choose which port of the large graph should be followed.

\begin{definition}[Zig-Zag Product~\cite{RVW2002}]
	\label{def:zig-zag}
	Let \(\Graph_1\) be an \((n_1,D_1,\lambda_1)\)-expander and let \(\Graph_2\) be a \((D_1,D_2,\lambda_2)\)-expander, where the vertex set of \(\Graph_2\) has size \(D_1\). The zig-zag product \(\Graph_1 \zzprod \Graph_2\) is the graph whose vertex set is \(\Vertices(\Graph_1)\times\Vertices(\Graph_2)\). A vertex \((v,c)\) should be read as the \(c\)-th cloud position over the base vertex \(v\). One step of \(\Graph_1 \zzprod \Graph_2\) is the three-part walk in which one first takes a step inside the \(\Graph_2\)-cloud, then follows the corresponding edge of \(\Graph_1\), and finally takes another step inside the destination \(\Graph_2\)-cloud. The resulting graph has \(n_1D_1\) vertices and degree \(D_2^2\).
\end{definition}

The words ``zig'' and ``zag'' refer to the two small steps inside the clouds, while the middle step is the movement in the large graph. The first small step randomizes the port through which the large graph is entered, and the second small step randomizes the cloud position after the large step has been taken. If \(\Graph_2\) expands well, these two short local corrections prevent the walk from remembering too much about a bad choice of port, which is why the product preserves expansion.

The quantitative statement below is the reason the zig-zag product can be iterated. The error term in the new spectral parameter is controlled by the old large-graph error and by the small cloud error, while the degree is reset to \(D_2^2\), and therefore the operation separates expansion from degree growth.

\begin{theorem}[Zig-Zag Expansion~\cite{RVW2002}]
	\label{thm:zig-zag}
	The zig-zag product \(\Graph_1 \zzprod \Graph_2\) is an \((n_1D_1,D_2^2,\lambda)\)-expander, where
	\begin{equation}
		\label{eq:zig-zag-spectral-bound}
		\lambda\leq \lambda_1+\lambda_2+\lambda_2^2\leq \lambda_1+2\lambda_2.
	\end{equation}
\end{theorem}

We stress that the key feature of \cref{thm:zig-zag} is that the degree of the product depends on \(D_2\), not on \(D_1\) and not on \(n_1\). Starting with a constant-size graph \(\Graph_2\) that already expands, and alternating zig-zag products with degree-amplifying operations such as graph squaring, one obtains explicit constant-degree expanders on arbitrarily many vertices. The construction is therefore not merely an existence proof; it is a recursive recipe whose neighbor function can be evaluated locally.

\begin{remark}
	The zig-zag product is the combinatorial engine behind Reingold's proof~\cite{Rei2008} that undirected \(s\)-\(t\) connectivity lies in \(\LOGSPACE\). The idea is to replace a graph \(\Graph\) by \(\Graph \zzprod H\) for a small expander \(H\), thereby reducing the degree while maintaining enough expansion, and then iterate this transformation until the walk can be simulated with logarithmic space. This application is conceptually important for the present note because it shows that explicit expansion is not only a pseudorandom object, but also a tool for controlling space-bounded computation.
\end{remark}

\subsection{The Replacement Product and Tensor Product}
\label{subsec:other-products}

Two other graph products arise repeatedly in cryptographic applications. The replacement product \(\Graph_1\circledR \Graph_2\) is a variant of the zig-zag product that reduces the degree from \(D_1\) to \(D_2+1\) while preserving expansion up to a constant-factor loss. Its importance is that it gives a flexible way to turn large-degree objects into bounded-degree objects, and it is one of the ingredients in the construction of lossless expanders~\cite{CRVW2002}. The tensor product, also called the categorical product, has vertex set \(\Vertices(\Graph_1)\times\Vertices(\Graph_2)\) and places an edge between \((u_1,u_2)\) and \((v_1,v_2)\) exactly when \(\{u_1,v_1\}\in\Edges(\Graph_1)\) and \(\{u_2,v_2\}\in\Edges(\Graph_2)\). Its spectrum consists of the products
\begin{equation}
	\label{eq:tensor-product-spectrum}
	\{\mu_i(\Graph_1)\mu_j(\Graph_2)\}_{i,j},
\end{equation}
which makes it useful when one wants to amplify spectral effects by multiplying eigenvalues.
